% Work/MAT — 1. ročník, hodina 62.
% Návaznost: 60 Definiční obor a obor hodnot; 61 Monotónnost, prostota,
% sudost, lichost. Příští hodina: Lineární funkce.
% Zdroj: ucebnice/02_Funkce_a_rovnice/Zaklady_funkci/Funkce.tex,
% oddíly Omezenost funkce, Maximum a minimum.
% Grafy a úlohy vlastní; znaménko funkce není předmětem této hodiny.
\documentclass[11pt,a4paper]{article}
% Spolecny styl pro ucitelske pripravy MAT (Work/MAT)
% Verze 1.1, 2026-10-23
\NeedsTeXFormat{LaTeX2e}
\ProvidesPackage{priprava}[2026/10/23 v1.1 Sablona ucitelskych priprav MAT]

% Zakladni typografie a cestina
\RequirePackage[T1]{fontenc}
\RequirePackage[utf8]{inputenc}
\RequirePackage[czech]{babel}
\RequirePackage{lmodern}
\RequirePackage{microtype}
\RequirePackage{amsmath,amssymb}
\RequirePackage{xcolor}
\RequirePackage{enumitem}
\RequirePackage[most]{tcolorbox}
\RequirePackage{tabularx,booktabs,array}
\RequirePackage{fancyhdr}
\RequirePackage{needspace}
\RequirePackage{geometry}
\RequirePackage[hidelinks]{hyperref}

\geometry{a4paper,top=17mm,bottom=18mm,left=23mm,right=23mm}
\definecolor{notegray}{gray}{0.39}
\setlength{\parindent}{0pt}
\setlength{\parskip}{2.5pt plus 1pt minus 0.5pt}
\setlist[itemize]{topsep=2pt,itemsep=1pt,parsep=0pt,leftmargin=1.45em}
\setlist[enumerate]{topsep=2pt,itemsep=2pt,parsep=0pt,leftmargin=1.65em}
\setcounter{secnumdepth}{0}
\renewcommand{\arraystretch}{1.12}

% Zahlavi: vlevo rocnik | globalni tema, vpravo cislo: nazev hodiny
% Zapati zustava zcela prazdne.
\pagestyle{fancy}
\fancyhf{}
\renewcommand{\headrulewidth}{0.35pt}
\renewcommand{\footrulewidth}{0pt}
\setlength{\headheight}{14pt}
\newcommand{\nastavPripravu}[4]{%
  \fancyhead[L]{\small #1\enspace\textbar\enspace #2}%
  \fancyhead[R]{\small #3: #4}%
}

% Minimalisticke cernobile boxy. Kazda definice zvlast.
\tcbset{commonlesson/.style={enhanced,breakable=false,colback=white,colframe=black,
  colbacktitle=white,coltitle=black,boxrule=0.45pt,
  arc=0.8mm,outer arc=0.8mm,left=3mm,right=3mm,top=1.6mm,bottom=1.6mm,
  toptitle=1mm,bottomtitle=0.7mm,fonttitle=\bfseries,
  attach boxed title to top left={xshift=2.5mm,yshift=-1.15mm},
  boxed title style={colback=white,colframe=white,boxrule=0pt,left=1mm,right=1mm}}}
\newtcolorbox{zadaniBox}[1][]{commonlesson,title={Zadání},#1}
\newtcolorbox{prikladBox}[1][]{commonlesson,title={Příklad},#1}
\newtcolorbox{definiceBox}[1][]{commonlesson,title={Definice},#1}

% Spolecny box pro zadani a jeho reseni, bez mezery mezi castmi.
% Pouziti:
%   \begin{prikladReseniBox}
%     Zadani ulohy.
%     \tcblower
%     \begin{align*} ... \end{align*}
%   \end{prikladReseniBox}
% Spodni cast ma jemny sedy podklad a delici linku; lze v ni sazet
% vice radku matematiky. Obe casti maji jen jeden spolecny ramecek.
\tcbset{solutionlesson/.style={
  commonlesson,
  enhanced jigsaw,
  breakable,
  colbacklower=black!3,
  segmentation style={solid,draw=black!25,line width=0.35pt},
  fontlower=\small\color{notegray},
  before lower={\noindent\textbf{Řešení:}\enspace},
  leftlower=3mm,rightlower=3mm
}}
\newtcolorbox{prikladReseniBox}[1][]{solutionlesson,title={Příklad},#1}
\newtcolorbox{zadaniReseniBox}[1][]{solutionlesson,title={Zadání},#1}

% Pomucky pro autora pripravy: sedive poznamky; \reseni pro starsi pripravy
\newcommand{\poznamka}[1]{\par\begingroup\small\color{notegray}\itshape
  \textnormal{\textbf{Poznámka pro učitele.}}\enspace #1\par\endgroup}
\newcommand{\reseni}[1]{\par\begingroup\small\color{notegray}
  \textbf{Očekávané řešení:}\enspace #1\par\endgroup}
\newcommand{\blok}[1]{\needspace{5\baselineskip}\subsection*{#1}}
\newcommand{\vyrok}[1]{\uv{#1}}

\usepackage{tikz}
\usetikzlibrary{arrows.meta}
\nastavPripravu{1. ročník}{Funkce}{62}{Omezenost a extrémy funkce}
\begin{document}
\poznamka{Jádro 45 minut: připomenutí oboru hodnot 5 min;
omezenost shora a zdola 10 min; tři porovnávací příklady 12 min;
maximum a minimum 9 min; samostatná práce a kontrola 7 min;
organizace 2 min. Náročnější případy a většina propojeného
procvičování jsou domácí úkol nebo rezerva.}

\blok{Kterých nejvyšších a nejnižších hodnot může funkce nabýt?}
V předchozích hodinách jsme se naučili číst $D(f)$ a $H(f)$
a posuzovat růst, pokles, sudost, lichost a prostotu.
Teď se zajímáme o to, \emph{zda jsou hodnoty nějak shora
či zdola ohraničené} a zda funkce dosahuje skutečně největší
nebo nejmenší hodnoty.

\begin{definiceBox}
Funkce $f$ je \textbf{shora omezená}, jestliže existuje
číslo $M\in\mathbb R$ takové, že pro každé $x\in D(f)$ platí
\[
f(x)\le M.
\]
Číslo $M$ nazýváme \textbf{horní mezí} funkce.
\end{definiceBox}
\begin{definiceBox}
Funkce $f$ je \textbf{zdola omezená}, jestliže existuje
číslo $m\in\mathbb R$ takové, že pro každé $x\in D(f)$ platí
\[
f(x)\ge m.
\]
Číslo $m$ nazýváme \textbf{dolní mezí} funkce.
\end{definiceBox}
\begin{definiceBox}
Funkce je \textbf{omezená}, jestliže je omezená současně
shora i zdola, tedy existují $m,M\in\mathbb R$ tak, že
\[
\forall x\in D(f):\qquad m\le f(x)\le M.
\]
\end{definiceBox}

\blok{Tři příklady omezenosti}
\begin{prikladReseniBox}[title={Příklad 1: pouze zdola omezená funkce}]
Urči omezenost funkce $f(x)=x^2$ na $\mathbb R$.
\tcblower
Pro každé $x\in\mathbb R$ je $x^2\ge0$; číslo $0$ je
proto dolní mezí. Funkce \textbf{je omezená zdola},
ale \textbf{není omezená shora}, protože pro dostatečně velké
$|x|$ překročí $x^2$ libovolné předem zvolené číslo $M$.
Obor hodnot je $H(f)=\langle0;+\infty)$.
\end{prikladReseniBox}

\begin{prikladReseniBox}[title={Příklad 2: pouze shora omezená funkce}]
Urči omezenost funkce $g(x)=-x^2$ na $\mathbb R$.
\tcblower
Protože $-x^2\le0$, číslo $0$ je horní mezí;
$g$ \textbf{je omezená shora}. Zdola omezená není,
neboť s rostoucí velikostí $|x|$ nabývá libovolně malých
(záporných) hodnot. $H(g)=(-\infty;0\rangle$.
\end{prikladReseniBox}

\begin{prikladReseniBox}[title={Příklad 3: omezená shora i zdola}]
Je dána $h(x)=x^2$ na $D(h)=\langle-2;1\rangle$.
Urči její omezenost a obor hodnot.
\tcblower
Na zadaném oboru je $0\le x^2\le4$:
nejmenší hodnotu $0$ dostaneme pro $x=0$, největší $4$
pro $x=-2$. Je tedy \textbf{omezená shora i zdola},
$H(h)=\langle0;4\rangle$. Samotný předpis je stejný
jako v příkladu 1, rozhodující je změna definičního oboru.
\end{prikladReseniBox}

\begin{center}
\begin{minipage}{.31\linewidth}\centering
\begin{tikzpicture}[x=.58cm,y=.43cm,>=Latex]
\draw[->] (-2.2,0)--(2.2,0) node[right] {$x$};
\draw[->] (0,-.5)--(0,4.5) node[above] {$y$};
\draw[thick,domain=-2:2,samples=65] plot (\x,{\x*\x});
\end{tikzpicture}\par\footnotesize $x^2$ na $\mathbb R$\\pouze zdola
\end{minipage}\hfill
\begin{minipage}{.31\linewidth}\centering
\begin{tikzpicture}[x=.58cm,y=.43cm,>=Latex]
\draw[->] (-2.2,0)--(2.2,0) node[right] {$x$};
\draw[->] (0,-4.5)--(0,.65) node[above] {$y$};
\draw[thick,domain=-2:2,samples=65] plot (\x,{-\x*\x});
\end{tikzpicture}\par\footnotesize $-x^2$ na $\mathbb R$\\pouze shora
\end{minipage}\hfill
\begin{minipage}{.31\linewidth}\centering
\begin{tikzpicture}[x=.58cm,y=.43cm,>=Latex]
\draw[->] (-2.2,0)--(2.2,0) node[right] {$x$};
\draw[->] (0,-.5)--(0,4.5) node[above] {$y$};
\draw[very thick,domain=-2:1,samples=65] plot (\x,{\x*\x});
\fill (-2,4) circle (2.8pt);\fill (1,1) circle (2.8pt);
\end{tikzpicture}\par\footnotesize $x^2$ na $\langle-2;1\rangle$\\oběma směry
\end{minipage}
\end{center}
\poznamka{Grafy jsou ilustrativní výřezy: první dvě funkce
pokračují mimo obrázek. Rozhodnutí o neomezenosti proto
nevyvozujeme jen z nakreslené části.}

\blok{Maximum a minimum: mez nemusí být dosažena}
\begin{definiceBox}
Funkce $f$ má \textbf{maximum}, existuje-li $x_0\in D(f)$ takové, že
\[
\forall x\in D(f):\qquad f(x)\le f(x_0).
\]
Číslo $f(x_0)$ je její největší funkční hodnota.
\end{definiceBox}
\begin{definiceBox}
Funkce $f$ má \textbf{minimum}, existuje-li $x_0\in D(f)$ takové, že
\[
\forall x\in D(f):\qquad f(x)\ge f(x_0).
\]
Číslo $f(x_0)$ je její nejmenší funkční hodnota.
\end{definiceBox}

\begin{prikladReseniBox}[title={Příklad 4: omezenost není totéž co maximum/minimum}]
Funkce $u(x)=x$ je zadána na $D(u)=(0;1)$.
Je omezená? Má maximum a minimum?
\tcblower
Platí $0<u(x)<1$, takže $u$ je omezená
shora i zdola. Nemá ale maximum ani minimum:
kdykoliv zvolíme $x\in(0;1)$, existují v tomto intervalu
jak větší, tak menší argumenty. Hodnoty $0$ a $1$
jsou mezemi, ale funkce jich nenabývá.
\end{prikladReseniBox}

\poznamka{Rozlišuj hodnotu extrému (např. maximum $4$)
a argument, kde nastává (např. $x=-2$).
Maximum se může dosahovat ve více bodech.}

\blok{Samostatná práce}
\begin{zadaniReseniBox}[title={Procvičování}]
\begin{enumerate}[label=\alph*)]
\item Je $p(x)=|x|$ na $\mathbb R$ omezená shora,
    zdola, nebo obojí? Má maximum či minimum?
\item Pro $q(x)=3-x$ na $\langle0;4\rangle$ urči
    $D(q)$, $H(q)$, maximum a minimum.
\item Funkce $r(x)=x$ na $(0;2\rangle$ je omezená?
    Má maximum nebo minimum?
\end{enumerate}
\tcblower
\begin{enumerate}[label=\alph*)]
\item Pouze zdola: $|x|\ge0$; minimum $0$ při $x=0$,
    shora neomezená, maximum nemá.
\item $D(q)=\langle0;4\rangle$, $H(q)=\langle-1;3\rangle$;
    maximum $3$ při $x=0$ a minimum $-1$ při $x=4$.
\item $0<r(x)\le2$, tedy omezená shora i zdola.
    Maximum $2$ nastává při $x=2$, minimum neexistuje.
\end{enumerate}
\end{zadaniReseniBox}

\blok{Domácí procvičení — souhrn hodin 59--62}
\begin{zadaniReseniBox}[title={Domácí úkol}]
\begin{enumerate}[label=\alph*)]
\item Funkce $f(x)=x^2$ je zadána na $\langle-2;1\rangle$.
    Urči $D(f)$, $H(f)$, intervaly růstu a poklesu,
    prostotu, sudost/lichost, maximum a minimum.
\item Pro $g(x)=x^3$ na $\mathbb R$ rozhodni:
    roste či klesá, je prostá, sudá či lichá,
    je shora nebo zdola omezená? Urči $D(g)$ a $H(g)$.
\item Pro $h(x)=-x^2$ na $\mathbb R$ určuj
    sudost/lichost, prostotu, omezenost a případné extrémy.
\item Funkce $p(x)=2x+1$ je definována na $\langle-1;2\rangle$.
    Urči $D(p)$, $H(p)$, monotonii, prostotu a extrémy.
    Je sudá nebo lichá?
\item Funkce $u(x)=x$ na $(0;1)$ je omezená shora i zdola.
    Vysvětli, proč přesto nemá maximum ani minimum.
\item Funkce $v$ je zadána tabulkou
    $v(-2)=4$, $v(-1)=1$, $v(1)=1$, $v(2)=4$
    na $D(v)=\{-2,-1,1,2\}$.
    Urči $H(v)$, sudost/lichost, prostotu a omezenost.
\item Z předpisu $w(x)=\dfrac{\sqrt{x-1}}{x-3}$
    urči přirozený definiční obor.
\end{enumerate}
\tcblower
\begin{enumerate}[label=\alph*)]
\item $D(f)=\langle-2;1\rangle$, $H(f)=\langle0;4\rangle$.
    Klesá na $\langle-2;0\rangle$, roste na $\langle0;1\rangle$.
    Není prostá ($f(-1)=f(1)$), ani sudá, ani lichá
    (nesouměrný obor). Maximum $4$ v $x=-2$,
    minimum $0$ v $x=0$.
\item $D(g)=H(g)=\mathbb R$. $g$ roste na $\mathbb R$,
    je prostá a lichá, nikoli sudá.
    Není omezená shora ani zdola a nemá maximum ani minimum.
\item $h$ je sudá, není prostá, je shora omezená číslem $0$,
    zdola omezená není. Má maximum $0$ při $x=0$,
    minimum nemá.
\item $D(p)=\langle-1;2\rangle$,
    $H(p)=\langle-1;5\rangle$; roste, je prostá,
    maximum $5$ má v $x=2$, minimum $-1$ v $x=-1$.
    Sudá ani lichá není, protože obor není souměrný vůči nule.
\item Pro každé $x\in(0;1)$ najdeme větší i menší
    přípustné číslo, takže $x$ nemůže být bodem
    maxima ani minima. Přesto $0<u(x)<1$.
\item $H(v)=\{1,4\}$; funkce je sudá
    ($v(-x)=v(x)$), není prostá
    ($v(-2)=v(2)$), je omezená shora i zdola
    (např. $1\le v(x)\le4$). Maximum $4$ a minimum $1$.
\item Pro odmocninu $x-1\ge0$ a pro jmenovatel $x\ne3$.
    Proto $D(w)=\langle1;3)\cup(3;+\infty)$.
\end{enumerate}
\end{zadaniReseniBox}
\end{document}
